\documentclass[11pt,a4paper]{article}
\usepackage[margin=2.15cm]{geometry}
\usepackage[T1]{fontenc}
\usepackage[utf8]{inputenc}
\usepackage{lmodern,microtype}
\usepackage{amsmath,amssymb,booktabs,tabularx,array}
\usepackage{graphicx,float,xcolor,hyperref}
\usepackage{enumitem}
\hypersetup{colorlinks=true,linkcolor=black,urlcolor=blue,citecolor=black}
\setlength{\parindent}{0pt}
\setlength{\parskip}{5pt}
\setlist[itemize]{topsep=2pt,itemsep=1pt,parsep=0pt}
\renewcommand{\arraystretch}{1.14}
\newcommand{\todo}[1]{\textcolor{red!65!black}{\textbf{[To verify: #1]}}}
\newcommand{\metric}[1]{\ensuremath{\mathrm{#1}}}
\title{Cross-Lingual Multi-Label Emotion Classification Using Large Language Models}
\author{Alex Smolders \and Felicia F\"alldin \and Thomas Platzer}
\date{November 2026}
\begin{document}
\maketitle

\begin{abstract}
Recognizing emotions from short dialogue is difficult because emotions can be implicit, several labels can apply to one utterance, and surrounding context is often missing. We study multilingual emotion classification using matched English--Swedish subtitle pairs from the XED dataset, with German considered as an additional language where completed results are available. We compare zero-shot and five-shot prompting with parameter-efficient fine-tuning (PEFT). The prompting experiments use Qwen2.5-3B-Instruct as the primary model and EuroLLM-9B-Instruct as a second model. The main study measures multi-label classification quality, prompt sensitivity, format compliance, and cross-lingual transfer. This manuscript retains preliminary results from the team's draft; these must be checked against the final exported metric files before submission. \todo{replace this sentence with a quantitative abstract after verification}
\end{abstract}

\section{Introduction and Motivation}
Emotions influence how people communicate, making it useful for language systems to recognize not only what is said but also the attitude conveyed by a statement. Emotion classification is relevant to applications such as conversational systems, online communication, and feedback analysis. However, predicting emotions from movie subtitles is not straightforward. A subtitle can be short, ambiguous, sarcastic, or dependent on previous dialogue that the model cannot see.

Our study addresses \emph{cross-lingual multi-label emotion classification}. The task is to identify every emotion expressed in a subtitle from Plutchik's eight categories: anger, anticipation, disgust, fear, joy, sadness, surprise, and trust. Unlike single-label classification, the model must decide which labels apply and how many labels to return. Moreover, good performance in English does not guarantee equally good performance in Swedish or German. Differences may arise from pretrained language coverage, translated expressions, and noise in the annotations, rather than language or culture alone.

We use the multilingual XED corpus \cite{ohman2020xed}, which contains human-annotated English examples and projected labels in additional languages. Our two central research questions are: (1) how do zero-shot prompting, five-shot prompting, and PEFT compare on multi-label emotion classification; and (2) how does performance change when the same task is evaluated in English and an aligned target language? We additionally examine sensitivity to prompt wording and the computational trade-offs of the methods.

\section{Background}
\subsection{Parameter-Efficient Fine-Tuning (PEFT)}
Parameter-efficient fine-tuning updates a small subset of a pretrained model rather than all of its parameters. The original weights are usually frozen, reducing the optimizer memory and the number of trained parameters. This is attractive for instruction-tuned language models that are too expensive to fully fine-tune on the available hardware.

\subsubsection{LoRA}
Low-Rank Adaptation (LoRA) \cite{hu2022lora} approximates a trainable weight update using two low-rank matrices. For a frozen projection matrix $W$, the adapted transformation is written as
\begin{equation}
W' = W + \frac{\alpha}{r}BA,
\end{equation}
where $r$ is the adapter rank and $\alpha$ controls the update scale. We target the attention projections \texttt{q\_proj} and \texttt{v\_proj}; the remaining pretrained weights are frozen. This preserves most of the pretrained model while allowing it to learn task-specific response patterns.

\subsubsection{QLoRA (Quantized LoRA)}
QLoRA combines LoRA adapters with low-bit storage of the frozen model weights \cite{dettmers2023qlora}. In an NF4-based implementation, the base model is stored in 4-bit precision while the training computation uses a suitable floating-point format. This can substantially reduce memory use. \textbf{Implementation note:} the original code archive called its method QLoRA but loaded 8-bit weights. The accompanying revised code changes the configured QLoRA path to 4-bit NF4. The results reported below must not be described as 4-bit QLoRA unless the corresponding models were actually retrained with this corrected implementation.

\subsubsection{Zero-Shot and Five-Shot Prompting}
Prompting adapts the model's behavior without changing its weights. In zero-shot prompting, the instruction and new subtitle are provided directly. Five-shot prompting additionally places five labeled training examples into the context, allowing the model to infer the required mapping from text to labels \cite{brown2020language}. Both approaches require a strict output contract: a JSON array containing only the eight allowed English emotion names. The same English-language instructions are used for English and Swedish input texts.

\section{Data and Task Setup}
\subsection{Data Gathering}
We use XED emotion annotations and sentence-aligned OpenSubtitles data extracted through OpusTools \cite{aulamo2020opus}. The pipeline downloads the target-language XED projection, extracts English--target one-to-one subtitle pairs from OpenSubtitles v2018, and matches them using normalized target-language text. The aligned English subtitle inherits the label set projected onto the matching target subtitle; it is not independently annotated in this pipeline. For the main study, the configured target language is Swedish (\texttt{sv}); German (\texttt{de}) can be prepared using the same pipeline with a separate configuration and outputs.

\subsection{Projected Label Status and Label Space}
XED includes eight emotion classes, with numerical IDs 1--8. The original neutral ID 0 is excluded from this eight-label task, and neutral mixed with another emotion is treated as invalid. In the resulting dataset, an utterance can have several emotion labels. Swedish reference labels are projections from aligned dialogue, so the reference data may contain alignment and projection errors. This matters particularly when a model's seemingly reasonable prediction differs from the recorded label set.

\begin{table}[H]
\centering
\caption{Mapping of XED IDs to emotion labels}
\begin{tabular}{cccccccc}
\toprule
1 & 2 & 3 & 4 & 5 & 6 & 7 & 8 \\
\midrule
anger & anticipation & disgust & fear & joy & sadness & surprise & trust \\
\bottomrule
\end{tabular}
\end{table}

\subsection{Normalization and Cleaning}
The pipeline applies Unicode NFKC normalization, trims whitespace, and collapses repeated spaces. It drops target texts with conflicting XED label sets and removes ambiguous target text when several OPUS alignment IDs could refer to the same utterance. Exact duplicate English sentences are reduced deterministically. These decisions reduce obvious duplication but may also discard valid data, so the retained sample counts should be reported from \texttt{outputs/report/sv\_dataset\_summary.csv}.

\subsection{Leakage Controls and Split Procedure}
After alignment and cleaning, the pipeline sorts records by alignment ID and applies a two-stage multi-label stratified split using seed 40. The intended ratios are 70\% training, 20\% development, and 10\% test. Each English--Swedish pair is placed in the same partition. The code checks that no exact normalized English or target sentence and no constructed alignment ID appears in more than one partition.

\textbf{Provenance limitation:} the current implementation constructs its ID from the OPUS source document and sentence ID. The course brief explicitly asks for the \emph{original XED alignment identifier}. The ZIP alone does not demonstrate that the OPUS-derived ID is equivalent. This mapping must be verified or the methodological limitation acknowledged before final submission. \todo{document source-ID verification or revise data provenance}

\begin{table}[H]
\centering\small
\caption{Dataset sizes to fill from the reproducible split summary; the source ZIP contains no processed dataset artifacts}
\label{tab:split}
\begin{tabular}{lrrr}
\toprule
Subset & Aligned pairs & Mean labels per subtitle & Unique alignment IDs \\
\midrule
Training & -- & -- & -- \\
Development & -- & -- & -- \\
Test & -- & -- & -- \\
\bottomrule
\end{tabular}
\end{table}

\subsection{Data formatting}
For each aligned pair, the pipeline creates the English and Swedish instruction--response records. The output is a JSON list of allowed emotion names. The bilingual PEFT condition joins the English and Swedish training instances, while monolingual conditions use only one language. The five-shot demonstrations are selected from training data only, using a fixed set of alignment IDs shared across both evaluation languages. The selected five-shot texts are rendered in the corresponding evaluation language.

\section{Methodology}
\subsection{PEFT}
We use Qwen2.5-3B-Instruct as the primary language model for both prompting and PEFT, making the comparison between adaptations interpretable. The three training conditions are English-only, Swedish-only, and bilingual. English-only PEFT is evaluated on English and Swedish; Swedish-only PEFT on Swedish; and bilingual PEFT on both English and Swedish. The selected adapter targets are the query and value projections, with rank 16, alpha 32, and dropout 0.05. The current configuration specifies learning rate $2\times 10^{-4}$, batch size 8, gradient accumulation 8, six epochs, and maximum sequence length 512. Training uses BF16 on compatible GPUs and seed 42. At the end of each epoch, evaluation loss is used to choose the best checkpoint.

\subsubsection{LoRA}
The PEFT configuration also supports standard LoRA without quantization. This is a useful comparison in principle, but the current ZIP does not include metric artifacts proving that a separate non-quantized LoRA experiment was completed. Therefore, this report treats LoRA as a supported method rather than claiming a verified quantitative LoRA-versus-QLoRA comparison. \todo{add verified LoRA runs if performed, or remove this comparison from final claims}

\subsubsection{Q-LoRA}
The revised implementation uses NF4 4-bit quantization with double quantization for the frozen base model, as specified in \texttt{configs/peft.yaml}. Only the LoRA adapter parameters are optimized. The adapter and the exact training YAML are stored after each run, allowing the model to be reloaded for inference. If previously saved adapters were trained using the original 8-bit loader, they must be labeled accordingly or retrained with the corrected code.

\subsection{Prompting with zero-shot and Five-shot}
Zero-shot and five-shot prompting are evaluated for both Qwen2.5-3B-Instruct and EuroLLM-9B-Instruct. Every prompt lists the eight permitted labels, explicitly permits multiple emotions, and requires a JSON array in English. The base prompt is tested alongside three wording variants: a paraphrase, a version without an explicit format example, and a shorter instruction. This permits controlled analysis of prompt sensitivity on the development set.

Five examples are selected from aligned training pairs with a fixed seed and label-coverage constraints; the same five IDs are used for both models and both input languages. When evaluating English, their English texts are shown, while Swedish evaluation uses the corresponding Swedish texts. We choose the final wording using mean development Micro F1 across the two evaluation languages for each model and shot condition, with sample Jaccard as a tie-breaker.

\subsection{Decoding settings}
We use the models' chat templates and deterministic greedy generation with \texttt{do\_sample=False}. The prompting configuration allows at most 64 generated tokens and uses batched inference. Inputs are padded on the left, and attention masks are passed to generation. A model's pretraining and decoding behavior may still differ slightly depending on the precise hardware and software versions; these should be recorded with the final experiment artifacts.

\subsection{Output parser}
We evaluate generated responses as predicted label sets. The shared parser requires one complete JSON array of strings, normalizes capitalization and whitespace, and removes duplicate labels. If the response is not a JSON array or includes a label outside the eight permitted categories, it is marked malformed and contributes an empty predicted set. The rate of strict whole-response JSON compliance is recorded separately from malformed rate. We do not map arbitrary synonyms such as ``anxiety'' to the label ``fear'', since this would change the task definition.

\subsection{Tools and Infrastructure}
The implementation uses Python, PyTorch, Hugging Face Transformers and PEFT, scikit-learn, pandas, NumPy, and Matplotlib. The experiments are intended for the University of Groningen H\'abr\'ok HPC cluster and its Slurm scheduler. Training and evaluation settings are documented in version-controlled YAML files, with dependencies defined in \texttt{pyproject.toml} and \texttt{uv.lock}. \todo{insert actual GPU, PyTorch, Transformers, PEFT, scikit-learn versions and training/evaluation runtimes from completed runs}

\section{Evaluation Protocol}
We compute all metrics from the same fixed-order, eight-label binary predictions. The primary metrics are Micro F1 and sample-averaged Jaccard; secondary metrics are Macro F1, micro precision and recall, and Hamming loss. We also report per-emotion precision, recall, F1, support, exact-match ratio, predicted label cardinality, empty-prediction rate, and malformed-output rate for prompting. Undefined precision and recall use \texttt{zero\_division=0}. For Jaccard, a pair of empty label sets contributes zero under the selected scikit-learn implementation; this choice is applied consistently across methods.

\begin{equation}
\metric{MicroF1}=\frac{2TP}{2TP+FP+FN},\qquad
\metric{Jaccard}_{i}=\frac{|G_i\cap P_i|}{|G_i\cup P_i|}.
\end{equation}
We report the arithmetic mean of sample Jaccard values over all examples. Development data are used for checkpoint and prompt selection, while the test set is reserved for the final comparison. Prompting and PEFT are rescored with the same parser and metric implementation. Single-seed results are reported without a repeated-run standard deviation, because the ZIP does not contain repeated independent seeds.

\subsection{Output Parsing and Instruction Adherence}
Malformed generation can reduce scores even when a model produces emotionally relevant words. We therefore distinguish empty but valid JSON arrays (\texttt{[]}) from outputs that violate the required format. When predictions contain invalid or unparseable text, the resulting label set is empty and the malformed rate is increased. Error analysis considers these failures separately from classification mistakes.

\subsection{Controlled comparison}
The source data, split seed, allowed label set, selected demonstration IDs, and inference settings are fixed within each language experiment. The same aligned evaluation records are used for models compared on the same input language. English and Swedish results are not treated as independent dataset samples when they are aligned versions of the same original line. A test comparison is only made after the prompt wording and PEFT checkpoint have been frozen on development data.

\section{Results}
\begin{table}[H]
\centering\small
\caption{Final primary-metric results template. Replace dashes with results from \texttt{sv\_test\_results.csv}}
\label{tab:primary}
\begin{tabular}{llcccc}
\toprule
& & \multicolumn{2}{c}{English input} & \multicolumn{2}{c}{Swedish input} \\
Model & Condition & Micro F1 & Sample Jaccard & Micro F1 & Sample Jaccard \\
\midrule
Qwen & Zero-shot & -- & -- & -- & -- \\
Qwen & Five-shot & -- & -- & -- & -- \\
EuroLLM & Zero-shot & -- & -- & -- & -- \\
EuroLLM & Five-shot & -- & -- & -- & -- \\
Qwen & English PEFT & -- & -- & -- & -- \\
Qwen & Swedish PEFT & n/a & n/a & -- & -- \\
Qwen & Bilingual PEFT & -- & -- & -- & -- \\
\bottomrule
\end{tabular}
\end{table}

Table~\ref{tab:primary} is the required primary-results structure. Unlike a single Macro F1 summary, it allows readers to evaluate the overall balance between recovered emotion labels and incorrect predictions. The two input languages are placed next to each other to emphasize that they are matched versions of the same subtitle pairs, not two unrelated corpora. The Swedish-only adapter is not required to be tested on English by the course matrix; a reverse-transfer experiment could be reported separately as an optional extension. The corresponding per-emotion table and F1 heatmaps should show whether weaker aggregate scores reflect consistently difficult labels or just one poorly predicted emotion.

\textbf{Status of the numbers below.} Table~\ref{tab:prelim} reproduces values already present in the team's incomplete PDF. These results are \emph{preliminary}: the provided ZIP contains no completed metric artifacts against which to verify them, and the original 8-bit/4-bit implementation discrepancy must be resolved before the final report. The primary required Micro F1 and sample Jaccard results should be inserted from the revised analysis outputs rather than inferred from Macro F1.

\begin{table}[H]
\centering\small
\caption{Preliminary Macro F1 values from the team's draft, Swedish experiment; test values with development values in parentheses, not independently verified}
\label{tab:prelim}
\begin{tabular}{llcc}
\toprule
Method & Training data & English & Swedish \\
\midrule
Zero-shot prompting & -- & 0.300 (0.302) & 0.238 (0.210) \\
Five-shot prompting & -- & 0.369 (0.390) & 0.337 (0.288) \\
\midrule
PEFT English-only & English & 0.412 (0.433) & 0.294 (0.329) \\
PEFT target-only & Swedish & 0.386 (0.414)$^\dagger$ & 0.312 (0.341) \\
PEFT bilingual & English + Swedish & 0.402 (0.429) & 0.300 (0.339) \\
\bottomrule
\end{tabular}
\vspace{2pt}
\begin{minipage}{0.95\linewidth}\footnotesize
$^\dagger$ The English value for target-only PEFT appears in the team's draft, but is outside the brief's minimum PEFT evaluation matrix; verify whether this extra condition was actually run before retaining it.
\end{minipage}
\end{table}

Five-shot prompting appears stronger than zero-shot prompting in the draft's reported Macro F1 for the primary model, especially on Swedish subtitles. However, the same table shows that PEFT rankings can differ by input language. These observations should be checked using the final primary metrics rather than treated as confirmed conclusions. The script \texttt{scripts/analyze\_results.py} writes a comprehensive results table including Micro F1, sample Jaccard, Macro F1, precision, recall, Hamming loss, and exact match.

\begin{table}[H]
\centering\small
\caption{Preliminary output behavior from the incomplete report; verify with final prediction files}
\begin{tabular}{lcccc}
\toprule
& \multicolumn{2}{c}{Empty predictions (\%)} & \multicolumn{2}{c}{Labels per subtitle}\\
Method & English & Swedish & English & Swedish \\
\midrule
Zero-shot & 18.4 & 24.0 & 0.94 & 0.80 \\
Five-shot & 2.4 & 5.6 & 1.50 & 1.28 \\
PEFT English-only & 0.0 & 0.0 & 1.18 & 1.03 \\
PEFT target-only & 0.0 & 0.0 & 1.29 & 1.23 \\
PEFT bilingual & 0.0 & 0.0 & 1.16 & 1.16 \\
Gold standard & -- & -- & 1.42 & 1.42 \\
\bottomrule
\end{tabular}
\end{table}

\IfFileExists{../outputs/report/test_micro_f1_sv_en.png}{
\begin{figure}[H]\centering\includegraphics[width=0.94\linewidth]{../outputs/report/test_micro_f1_sv_en.png}\caption{Final English-input test Micro F1 comparison, generated from saved results}\end{figure}
}{}
\todo{insert verified main results table, prompt sensitivity figure, per-emotion F1 chart, model runtime, and German results if applicable}

\section{Analysis and Discussion}
The preliminary results suggest that providing five demonstrations may help a model emit more emotion labels and follow the task structure. Yet improving label coverage can also increase false positives. This trade-off should be evaluated using both micro precision and recall, not only Macro F1. The completed report should analyze differences between English and Swedish on the same aligned examples, as well as the relative strengths of prompting and PEFT.

\textbf{Cross-lingual transfer.} English-only PEFT evaluated on Swedish tests whether emotion patterns learned in English generalize without Swedish fine-tuning. Target-only and bilingual models offer two alternative adaptation strategies. Differences between these conditions may reflect the amount and type of training data, model pretraining, or properties of the label projection; they do not, by themselves, establish cultural differences.

\textbf{Prompt sensitivity.} The three alternative instructions provide a development-only check of how stable results are to wording. The report should include the mean Micro F1 and sample Jaccard for each variant and identify the selected wording. A large change across variants indicates that model performance depends on the phrasing of instructions rather than only the underlying ability to recognize emotions.

\textbf{Qualitative analysis.} We compare predictions at the same alignment ID when the prompted primary model is right but the PEFT model is wrong, and vice versa. The analysis script writes these changes to \texttt{outputs/report/sv\_test\_qualitative.csv}. We should choose concrete cases showing missing implicit emotions, extra predicted labels, incorrect labels, and malformed responses, and consider whether the gold labels themselves may have been affected by projection or incomplete dialogue context. \todo{insert 3--5 verified examples and discuss them individually}

\textbf{Computational cost.} Fine-tuning requires additional GPU time and saved adapters, whereas prompting avoids training but may require longer contexts in five-shot evaluation. The comparison should report measured training runtime, prediction runtime, peak GPU memory if recorded, and hardware. The original archive contains no comprehensive runtime log. The revised runner records prompting inference seconds and peak GPU allocation; new PEFT runs also record training seconds, peak GPU allocation, and inference seconds in their metadata. Slurm's \texttt{sacct} can provide an independent elapsed-time measurement. The analysis exports \texttt{sv\_test\_runtime.csv} with recorded values and leaves unavailable fields blank.

\subsection*{Interpreting different kinds of errors}
A central difficulty with a multi-label task is that a single prediction can contain both a correct label and an incorrect one. For instance, if the reference set contains fear and surprise but the model predicts only fear, the answer has perfect precision for this example but incomplete recall. By contrast, predicting fear, surprise, and anger covers both reference labels but introduces a false positive. Exact-match accuracy treats both predictions as incorrect, whereas Micro F1 and Jaccard distinguish their partial correctness. This is why we use the two primary metrics together rather than relying on exact match alone.

A second issue is the frequency of empty predictions. An instruction-tuned model may be reluctant to infer an emotion from a sentence without clear affective words, which can yield reasonable-looking but overly conservative outputs. In that case, its precision might be higher than its recall. Conversely, a detailed prompt that encourages implied emotions might improve recall at the cost of false positives. The recorded average number of predicted labels per subtitle therefore helps explain changes in classification quality. The gold-standard cardinality should be calculated directly from each matched test partition so that comparisons are not affected by differences in retained examples.

A third source of error is output formatting. The model can produce the right label in an invalid representation, such as an unquoted string inside brackets. Our scoring pipeline treats responses without parseable JSON as malformed and assigns an empty predicted set. This is deliberately stricter than searching arbitrary text for emotion words: broad recovery rules could hide instruction-following weaknesses and can accidentally treat explanations as classifications. We therefore report both the malformed-output rate and the stricter whole-response JSON compliance rate. Both should be inspected before attributing poor F1 entirely to misunderstanding of emotional content.

\subsection*{Validity and limits of comparison}
All prompted and fine-tuned models are evaluated with the same label order and the same held-out items for a given input language. Nevertheless, the amount of supervision differs. Five-shot prompting adds only five examples to the inference context and does not update model parameters, while bilingual fine-tuning uses both versions of the training corpus and therefore has more training examples per epoch than a monolingual condition. Better bilingual performance would not automatically prove that combining languages is better than using more data; a length-matched downsampling experiment would be needed to separate these effects.

The projected-label nature of XED also limits the interpretation of individual errors. The assigned-language sentence may not express precisely the same attitude as its aligned English version, and the source dialogue may rely on earlier lines not included in the input. If the model disagrees with a projected label, this can reflect an incorrect prediction, a translation mismatch, or a limitation of the reference annotation. We do not assume that the model is right merely because a label appears unintuitive. Instead, the qualitative analysis should present the actual sentence, reference set, predicted sets, and any relevant alignment uncertainties, while avoiding unsupported conclusions about cultural differences.

Finally, the current work is primarily a controlled, single-seed study, not a large hyperparameter search. Exact replication depends on frozen data partitions, prompt texts, five-shot IDs, model revisions, CUDA-compatible libraries, and adapter checkpoints. The final report should give the exact configuration and code commit associated with each headline result. Computational efficiency claims require measured runtime or memory rather than inferred advantages from parameter count alone.

\subsection{Why no chain-of-thought}
We do not ask the language models to output chain-of-thought explanations. The required output is a short JSON list, and free-form explanations would complicate output parsing and instruction-adherence measurements. Additional generated reasoning could also increase token use and produce unsupported interpretations of a context-free subtitle. This design choice keeps prompting and fine-tuning focused on the same explicit classification target; it does not mean that the models are unable to perform internal reasoning.

\section{Conclusion and Future work}
This project compares prompting and parameter-efficient adaptation for multi-label emotion classification in matched English and Swedish subtitles. The available draft suggests that a few demonstrations can help prompting, but the final conclusion must be based on reproducible Micro F1 and sample Jaccard results from the complete experiment matrix. The analysis should consider performance, malformed outputs, per-emotion differences, and computational cost together. Future work could test repeated training seeds, alternative adapter ranks or learning rates, verified original XED alignment metadata, and a carefully documented extension that adds surrounding dialogue context. \todo{replace with final evidence-based conclusion after all checks}

\begin{thebibliography}{9}
\bibitem{ohman2020xed} Emily \"Ohman, Marc P\`amies, Kaisla Kajava, and J\"org Tiedemann. XED: A multilingual dataset for sentiment analysis and emotion detection. \emph{Proceedings of COLING}, 2020.
\bibitem{aulamo2020opus} Mikko Aulamo, Umut Sulubacak, Sami Virpioja, and J\"org Tiedemann. OpusTools and parallel corpus diagnostics. \emph{LREC}, 2020.
\bibitem{brown2020language} Tom B. Brown et al. Language Models are Few-Shot Learners. \emph{NeurIPS}, 2020.
\bibitem{hu2022lora} Edward J. Hu et al. LoRA: Low-Rank Adaptation of Large Language Models. \emph{ICLR}, 2022.
\bibitem{dettmers2023qlora} Tim Dettmers, Artidoro Pagnoni, Ari Holtzman, and Luke Zettlemoyer. QLoRA: Efficient Finetuning of Quantized LLMs. \emph{NeurIPS}, 2023.
\end{thebibliography}

\appendix
\section{Prompt templates and reproducibility}
The exact system instructions and their wording variants are defined in \texttt{src/llm\_project/prompting.py}. The selected demonstration IDs, model identifiers and generation parameters are versioned in \texttt{configs/prompting.yaml}. The data split configuration is in \texttt{configs/data.yaml} and the PEFT setup in \texttt{configs/peft.yaml}. The README and RUNNING\_GUIDE reproduce the command sequence. The complete prompt cards are produced as Markdown files alongside prompting metrics.
\end{document}
